\section{Sharp lift ambiguity and the degree-two boundary}\label{sec:lift-boundary}
The bounds in this section use independent uniform matchings above each
base edge. Statistical claims hide the planted fibers by a uniform
relabeling before revealing the graph.

\subsection{Complete bases attain the structural failure rate}
\begin{theorem}[Complete-base leading constant]\label{thm:complete-lift-sharp}
Fix $d\ge3$, let $k=d+1$, and let $G$ be the random $N$-lift of
$K_k$. For any fixed $K\ge k$, let $\mathcal A_K$ be the event that
an equitable partition into at most $K$ classes does not coarsen the plant.
Then, with $\beta=d-2$ and $c_d=\binom{d+1}{2}$,
\[
 \Prob(\mathcal A_K)=c_dN^{-\beta}+o(N^{-\beta}).
\]
If the known base is $K_k$ and the balanced plant is hidden uniformly,
the optimal recovery Bayes error is
$(c_d/2)N^{-\beta}+o(N^{-\beta})$. Both leading asymptotics persist
after conditioning the total lift connected.
\end{theorem}
\begin{proof}
For $u\in V_i,v\in V_j$, $i\ne j$, let $I_{uv}$ assert adjacency
and equality of their neighbors in every other old fiber. Independence
of the base-edge matchings gives $\Prob(I_{uv})=N^{-d}$.
On this event the closed neighborhoods of $u,v$ agree, so transposing
them is a graph automorphism. Swapping them between their fibers gives
a different balanced cover partition for $N\ge2$. If $X$ counts these
unordered cross-fiber twin pairs, then exactly
\[
 \E X=\binom k2N^{-(d-2)}.
\]
One pair corresponds to a colored motif $F$ consisting of its two
centers, their edge, and $d-1$ common neighbors, one in each remaining
fiber. Thus $v(F)=d+1$, $e(F)=2d-1$, and its excess is $\beta$.

For two distinct feasible motifs with common graph $J$, include all
common vertices, even isolated ones, and put $m=v(J),h=e(J)$.
A fixed overlap pattern has $O_d(N^{2k-m})$ placements and probability
$O_d(N^{-2(2d-1)+h})$, by the exact bounded partial-permutation
probabilities $\prod_{ij}(N)_{r_{ij}}^{-1}$.
Moreover $h\le m$. If the centers coincide, matching injectivity
forces the whole motif to coincide, which is excluded. Otherwise
one center of the first motif is absent or has degree at most two in
the second. Removing it from the first motif leaves a star. Hence
the common graph is a forest with at most one added cycle, proving
$h\le m$. Summing the finitely many overlap types yields
\[
 \E[X(X-1)]=O_d(N^{-2\beta}),\qquad
 \Prob(X\ge1)=c_dN^{-\beta}+O_d(N^{-2\beta}),\quad
 \Prob(X\ge2)=O_d(N^{-2\beta}).
\]

We refine the rigidity proof of Section~\ref{sec:unknown-frame}.
Its macroscopic alternative event has exponentially small probability.
For small bad supports, \eqref{eq:unknown-support} now sums from
size three, giving $O_{d,K}(N^{-3\beta/2})=o(N^{-\beta})$.
A singleton is impossible. A bad support of size two is exactly one
of the twin pairs above. Two vertices of the same old fiber cannot
share a neighbor, by matching injectivity. For vertices in different
fibers that are not adjacent, the neighbor of the first in the second
fiber is exterior and has exactly one support neighbor. If they are
adjacent, badness says precisely that all other neighbors coincide.

Outside the macroscopic event, every candidate $Y$ is within $\eta n$
rows of a point map $JT$. Take $\eta<1/[2k(d+1)]$. Every old fiber
then contains an unchanged closed neighborhood, proving $QT=TR$.
Its nonzero difference $Y-JT$ has a bad row support. Outside the
size-at-least-three event it is a twin pair. Conversely every twin
pair supplies a noncoarsening candidate. These inclusions prove the
ambiguity asymptotic, with error $O_{d,K}(N^{-3\beta/2})$.

For the Bayes assertion, every ordered balanced $K_k$-cover partition
has constant likelihood $(N!)^{-\binom k2}$ on compatible labeled
graphs. Each unordered one has $k!$ orderings, so the posterior is
uniform on its $Z(G)$ compatible balanced cover partitions.
Near a point map, a balanced $k$-class candidate must use an onto
map: an omitted label would force at least $N$ changes. Thus that map
is a permutation. After naming it identically, a size-two difference
preserving all class sizes can only swap the two twin vertices.
Outside the previous exceptional events and $\{X\ge2\}$, $Z=1+X$.
Consequently $\Prob(Z\ge3)=o(N^{-\beta})$ and optimal error
$\E(1-1/Z)\sim \frac12c_dN^{-\beta}$.

Finally every component of a cover meets each fiber in the same
number $m$. A disconnected lift has one with $1\le m\le N/2$.
For specified $m$-subsets in every fiber, all matchings preserve them
with probability $\binom Nm^{-\binom k2}$. Union bounding gives
\[
 \Prob(G\text{ disconnected})
 \le\sum_{m=1}^{\lfloor N/2\rfloor}\binom Nm^{-r_0}
 =O_d(N^{-r_0}),\qquad r_0=\binom k2-k=k(d-2)/2>\beta.
\]
For the displayed order, separate $m=1$ and use
$\binom Nm\ge\binom N2$ for $2\le m\le N/2$.
Conditioning connected therefore preserves both leading asymptotics,
and does not change the uniform posterior on a fixed connected graph.
\end{proof}

The Bayes constant here assumes the base is supplied. It also gives an
information lower bound for any unknown-base class containing this
ensemble. A posterior over unknown bases would require its own prior.
The complete base shows that the universal structural exponent $d-2$
cannot generally be increased.

\subsection{Exact failure of recovery for cycle bases}
\begin{theorem}[Known cycle, hidden fibers]\label{thm:cycle-boundary}
Fix $k\ge3$ and consider the unconditioned random $N$-lift of the
known simple cycle $C_k$, followed by a uniform relabeling. If $R(G)$
is its component count, the posterior on compatible unordered fiber
partitions is uniform on $(2k)^{R(G)-1}$ possibilities. Optimal Bayes
success, and optimal randomized minimax success, are exactly
\begin{equation}\label{eq:cycle-success}
 P_{k,N}=\frac{2k\,\Gamma(N+1/(2k))}
 {\Gamma(1/(2k))\Gamma(N+1)}
 \sim\frac{2k}{\Gamma(1/(2k))}N^{-1+1/(2k)}.
\end{equation}
Thus success tends to zero without any computational restriction.
Connectedness has probability $1/N$; conditional on connectedness,
the supplied-base unordered fiber partition is uniquely recoverable.
\end{theorem}
\begin{proof}
Transport names along a spanning path of the base, making its first
$k-1$ matchings identities. The last matching is a uniform permutation
$\pi\in S_N$. Its cycles of lengths $\ell_1,\ldots,\ell_R$ give graph
components $C_{k\ell_1},\ldots,C_{k\ell_R}$.

A component $C_{k\ell}$ has exactly $2k$ locally bijective maps to
the named base: choose one vertex's image in $k$ ways and a direction
in two ways. Local bijectivity forces that direction thereafter.
Every such map puts $\ell$ vertices in each fiber. Thus all components
give $(2k)^R$ balanced named maps. Two maps have the same unordered
fibers exactly when they differ by a single global automorphism of
$C_k$. Indeed the induced permutation of fiber names must preserve all
base edges; conversely any base automorphism works. This action is
free and has order $2k$. Repeated component lengths create no extra
quotient, since observed vertices are labeled. The partition count is
therefore $(2k)^{R-1}$.

There are $n!/(N!)^k$ balanced named maps. The prior probability of one
is $(N!)^k/n!$; its likelihood for a compatible graph is $(N!)^{-k}$.
Every compatible map thus has joint probability $1/n!$. Summing over
its $2k$ global-name variants gives the same joint probability to each
compatible unordered partition, proving the posterior claim.

For the cycle count $R_N$ of a uniform permutation, inserting the last
element as a singleton or after one of the other $N-1$ elements gives
\[
 \E t^{R_N}=\frac{t(t+1)\cdots(t+N-1)}{N!}
 =\frac{\Gamma(N+t)}{\Gamma(t)\Gamma(N+1)}.
\]
Take $t=1/(2k)$ and multiply by $2k$ for \eqref{eq:cycle-success};
Stirling gives its asymptotic. A permutation is one cycle with
probability $1/N$, and at $R=1$ the partition count is one.

Finally the uniform hidden prior bounds any worst-case success from
above by $P_{k,N}$. Sampling uniformly from the compatible partitions
attains that success for every fixed hidden named map: fix a canonical
phase and direction on one component, and choose the $2k$ possibilities
independently on each other component. Thus it is randomized minimax
optimal. The hidden-prior or minimax formulation is essential; a fixed
publicly named plant can be hardcoded by an estimator.
\end{proof}

Connected conditioning is a different degree-two model. If the cycle
base size is unknown, a connected cycle can cover other cycle bases;
no identification among those bases is asserted by the last theorem.

